\section{Conclusion}\label{sec:conclusion}

The same support geometry controls the exact and approximate problems.
On an \((m+1)\)-party marginal it exposes every local Weyl axis, which
turns arbitrary product-unitary intertwiners into product Cliffords.  The
transition maps between the resulting local Weyl frames organize the
residual freedom.  After one party is reinterpreted as a logical input, the
support geometry becomes
three-region correctability.  Cleaning and Weyl--Fourier concentration
then recover the same discrete Clifford frame without paying the
exponentially small amplitude of a single marginal.

The exact classification is finite and constructive.  A half-set direct sum
reduces the data to the dimensionally optimal \(m\) minimum supports and gives
fixed-label recognition with an explicit field-operation bound.  In prime
dimension, a quadratic form on four variables reduces decision and witness
construction to polynomial cost in \(m\) and \(\log q\), with an expected
bound for construction.  Over
prime fields, the common centralizer of the fundamental cycles is the
intrinsic algebra of block-diagonal one-party label endomorphisms preserving
the stabilizer.  Its five possible types determine both the symmetry group
and the additional code structure.  For four parties it is the full matrix algebra; for six it is necessarily
nonscalar.  The same recognition procedure decides transversal
conversion between the associated encoders and stochastic conversion
between the states, since every product operator between two such states
is Clifford up to scalars.  The same \(m\) minimum supports carry the
reductions that determine the state among all density operators, at the
smallest locality any parent Hamiltonian can have.  Testing both complementary
families raises the verification gap to \((m+1)/(2m)\).  This turns the
marginal constraints into a fidelity certificate and, under a product-unitary
preparation promise, an observable entry condition for robust rigidity.

The quantitative conclusion is uniform only when all varying parameters
are displayed:
\[
 R^{\mathrm{clean}}_{n,q}
 =\Theta\!\left(\min\{q^{-1/2},n^{-1/2}\}\right),
 \qquad q=p^e.
\]
The length restriction \(n\leq2(q^2-1)\) excludes unbounded fixed-dimension
families.  On generalized and extended
Reed--Solomon constructions, where \(n\leq q+1\), the scale is
\(\Theta(q^{-1/2})\) over both prime and extension fields.  Proposition~\ref{prop:spectral-chart-obstruction} shows that the dimension
exponent is necessary when both the spectral-spread condition and collective
coefficient are retained.  The matching bounds in
Remark~\ref{rem:radius-sandwich} give order-sharpness whenever \(n\leq Cq\);
they need not match in regimes where \(n/q\) grows.

The robust transition-map result identifies the remaining bottleneck.
Linear symplectic compatibility becomes exact for
\(\varepsilon<\sqrt2/68\), independently of party count and local dimension, but commutators discard the affine Weyl phases.  They
therefore determine the stabilizer label group and not its character.  The
logical action of any chosen encoder can nevertheless be rounded within
\(8\varepsilon\) to an exactly transversally realizable Clifford, because a
stabilizer cancels the Pauli correction on the logical leg.  A
party-count-independent global entry theorem would need a new way to
control that character discrepancy, or an equivalent collective branch
selection principle, before the product-Pauli correction is known to be
locally small.
